% Research-direction note: strategies for proving the global tripartite
% Turan-filling reduction for P = K_4^dagger and J = K_3 cup_{K_2} K_4.

\documentclass[11pt]{amsart}

\usepackage{amsmath,amssymb,amsthm,mathtools}
\usepackage[margin=1.1in]{geometry}
\usepackage{enumitem}
\usepackage{hyperref}
\usepackage{xcolor}

\hypersetup{
  colorlinks=true,
  linkcolor=blue!50!black,
  citecolor=blue!50!black,
  urlcolor=blue!50!black
}

\newtheorem{theorem}{Theorem}[section]
\newtheorem{proposition}[theorem]{Proposition}
\newtheorem{lemma}[theorem]{Lemma}
\newtheorem{conjecture}[theorem]{Conjecture}
\newtheorem{problem}[theorem]{Problem}
\newtheorem{corollary}[theorem]{Corollary}

\theoremstyle{definition}
\newtheorem{definition}[theorem]{Definition}
\newtheorem{remark}[theorem]{Remark}

\newcommand{\graphon}{\mathcal W}
\newcommand{\E}{\mathbb E}
\newcommand{\R}{\mathbb R}
\newcommand{\dd}{\,d}
\newcommand{\Hom}{\operatorname{Hom}}

\title[Strategies for the global tripartite-filling reduction]%
{Strategies for the global tripartite Turán-filling reduction:\\
$f_P(x)=\Psi_P(x)$ and $f_J(x)=\Psi_J(x)$ on $[2/3, 3/4]$}
\author{Companion strategy note}
\date{\today}

\begin{document}
\maketitle

\begin{abstract}
This note records the research-level structure of the open global reduction
\[
   f_P(x) = \Psi_P(x),\qquad f_J(x) = \Psi_J(x)\qquad\text{on}\ [2/3,3/4],
\]
where $P = K_4^{\dagger}$, $J = K_3\cup_{K_2}K_4$, and $\Psi_P,\Psi_J$ are the
explicit one-dimensional reduced curves of §7 of the progress note.  We
formulate the key intermediate claims, identify the precise asymmetries that
make this problem harder than the clique case, list three independent attack
strategies, and isolate concrete lemmas to be proved or refuted next.

The note is deliberately not a finished proof; it is a research roadmap.
\end{abstract}

\tableofcontents

\section{The conjecture and a useful integral identity}

For brevity write $P = K_4^{\dagger}$ ($K_4$ with one pendant edge attached to
a fixed $K_4$-vertex) and $J = K_3\cup_{K_2}K_4$ ($K_3$ and $K_4$ glued
along a common edge).

\begin{conjecture}[Global tripartite-filling reduction]
\label{conj:global-reduction}
For every $x\in[2/3,3/4]$,
\[
   f_P(x) = \min\{t(P,W): t(K_2,W)=x\} = \Psi_P(x),
\]
where $\Psi_P$ is the closed-form reduced curve of §7.4 of the progress note,
and the minimum is achieved by a graphon in the tripartite Turán-filling
family
\[
   \mathcal T = \biggl\{\ W\ \biggm|\ \exists [0,1] = A\sqcup B\sqcup C, |B|=|C|,
   \begin{aligned}
   &W=1\ \text{on cross-block cells},\\
   &W=0\ \text{on within-}B\text{ and within-}C,\\
   &W|_{A\times A}\ \text{triangle-free}
   \end{aligned}\ \biggr\}.
\]
The analogous statement holds for $J$ and $\Psi_J$.
\end{conjecture}

\subsection{A rooted-degree integral identity}

Both $t(P,\cdot)$ and $t(J,\cdot)$ can be written as integrals of a "rooted
weight" against a "local quantity" at a distinguished vertex.  This is the key
reason their minimisers are not the pure clique-density minimisers, and it is
the lever on which any structural reduction must act.

For $P = K_4^{\dagger}$ with the pendant attached to vertex $4$ of the $K_4$,
\begin{equation}
\label{eq:P-integral}
   t(P,W)
   = \int_{[0,1]} d_W(x)\,\kappa_4(x)\dd x,
\end{equation}
where
\[
   d_W(x) := \int W(x,y)\dd y,\qquad
   \kappa_4(x) := \int \prod_{ij\in E(K_4),\ i,j\ne 4}\hspace{-2em} W(x_i,x_j)\hspace{-0.5em}
   \prod_{i\ne 4} W(x,x_i)\dd x_1\dots\widehat{\dd x_4}\dots\dd x_4
\]
is the \emph{rooted $K_4$-density at $x$} (the probability that three uniformly
random partners of $x$ form a $K_3$ together with $x$).  Symmetrically
$t(K_4,W) = \int\kappa_4(x)\dd x = \E[\kappa_4(X)]$.

Likewise, $J$ can be written as an integral over the shared edge:
\begin{equation}
\label{eq:J-integral}
   t(J,W) = \int_{[0,1]^2}\hspace{-1em} W(x,y)\,\tau_3(x,y)\,\tau_4(x,y)\dd x\dd y,
\end{equation}
where $\tau_3(x,y)$ is the rooted $K_3$-density at the edge $\{x,y\}$ (probability that a
random third vertex closes a triangle with $x,y$) and $\tau_4(x,y)$ the rooted
$K_4$-density at $\{x,y\}$ (probability that two random additional vertices
form a $K_4$ with $x,y$).

\subsection{Why \protect{$f_P\ne f_{K_4}$ in general}}

Cauchy-Schwarz applied to~\eqref{eq:P-integral} at fixed
$\E[d_W]=2x$ and fixed $\E[\kappa_4]=t(K_4,W)$ gives only an upper bound on
$t(P,W)$.  The lower bound on $t(P,W)$ at fixed $t(K_2,W)$ is achieved when
$d_W$ and $\kappa_4$ are \emph{anti-correlated} (high degree at low rooted
clique density, and vice versa), and the magnitude of the correlation is what
the tripartite-filling family exploits.

In the Lovász--Simonovits clique template, $d_W$ and $\kappa_4$ are essentially
constant across each part, and the few distinct values are positively
correlated (the small "singleton" part has both highest $\kappa_4$ and highest
$d_W$).  This is exactly the wrong correlation for minimising $t(P,W)$, which
is why the LS template is not optimal for $P$.

In the tripartite-filling family at our solution:
\begin{itemize}[itemsep=2pt]
\item Vertices in $A$ have degree $d = 1-\alpha+\alpha q$ and rooted-$K_4$
density largely determined by their connection pattern in $A$.  Because $W|_{A\times A}$
is triangle-free, rooted-$K_3$ counts inside $A$ vanish, suppressing the
rooted $K_4$-density at $A$-vertices.
\item Vertices in $B$ (and $C$) have degree $(1+\alpha)/2$ and rooted-$K_4$
density given by all 4-cliques whose root sits in $B$.
\end{itemize}
The construction is the simplest way to lower $\kappa_4$ on the high-degree
vertices and raise it on the low-degree ones.

\section{Plan A: stability + symmetrisation}

This route is the most natural one and is closest to the clique case
(Reiher 2016).  The two non-trivial steps are:

\begin{problem}[Step A1: stability of a 3-blowup skeleton]
\label{prob:A1}
Let $W$ be a graphon with $t(K_2,W)=x\in[2/3,3/4]$ minimising
$t(P,W)$ (resp.\ $t(J,W)$).  Show that $W$ has a "3-blowup skeleton":
$[0,1] = A\sqcup B\sqcup C$ measurable, with
\begin{enumerate}[label=(\arabic*),itemsep=2pt]
\item $W=1$ a.e.\ on the three cross-block products $A\times B, A\times C,
B\times C$;
\item $W=0$ a.e.\ on the within-block products $B\times B, C\times C$.
\end{enumerate}
The internal structure inside $A$ is unconstrained at this point.
\end{problem}

\begin{problem}[Step A2: symmetrisation to balanced parts]
\label{prob:A2}
With the skeleton above fixed, show that one may further assume $|B|=|C|$
without increasing $t(P,W)$ (resp.\ $t(J,W)$), and that $W|_{A\times A}$ may
be replaced by a $q$-regular triangle-free filling with the same edge density.
\end{problem}

If Steps A1 and A2 both hold, then $f_P(x) \le \Psi_P^{\text{filling-only}}(x)$
where the right-hand side is the minimum over the tripartite-filling family.
Combined with the lower bound $f_P\ge \Psi_P^{\text{filling-only}}$ inside the
family (i.e.\ the reduced problem of §7), this gives $f_P=\Psi_P$ on $[2/3,3/4]$.

\subsection{Step A1: clique-stability route}

Step~A1 is closely related to clique-density supersaturation.  Define
\[
   d_+(W) := \mathrm{vol}\bigl(\{(x,y) : W(x,y)>0,\ \text{but $W$ is "far" from a 3-blowup}\}\bigr).
\]
A clique-stability statement of Pikhurko--Razborov type bounds $d_+$ by a
function of the supersaturation gap $t(K_4,W) - g_{4,3}(s_*)$ where
$g_{4,3}$ is the LS curve.  Concretely, the available tool is:

\begin{lemma}[Stability for $K_4$ supersaturation on the first Turán interval, conjectural form]
\label{lem:K4-stability}
There is a constant $C(x)>0$ such that for every $W$ with $t(K_2,W)=x\in(2/3,3/4)$,
\[
   \delta_\square\bigl(W,\ \mathcal{T}_{3}^{1}(x)\bigr)
   \le C(x)\cdot \bigl(t(K_4,W) - g_{4,3}(s_x)\bigr)^{1/2},
\]
where $\delta_\square$ is the cut metric, $\mathcal{T}_3^1(x)$ denotes the set of
3-blowups (no constraint on the internal structure inside $A$) of total edge
density $x$, and $g_{4,3}(s_x)$ is the LS $K_4$-density at $x$ with the 4-partite
template (computed in §7 of the progress note).
\end{lemma}

\begin{remark}[Status of Lemma~\ref{lem:K4-stability}]
A stability statement of this strength for the LS 4-partite template at
$K_4$-density is not directly stated in the literature.  Pikhurko--Razborov
type results cover the equality case at the
clique-density level; the cut-metric stability above would require either
adapting the standard regularity/decomposition argument used for triangle
supersaturation to the 4-partite case, or invoking a more recent stability
result for clique-density.  It is the most technical ingredient.
\end{remark}

\subsection{From skeleton to balanced parts}

If $W$ has a 3-blowup skeleton, the contribution to $t(P,W)$ splits cleanly
into "cross-edges only" (the LS-template-like contribution) plus a "filling
inside $A$" contribution.  The filling-only contribution is the quantity that
the §7 analysis controls.  Step A2 is then a "balancing the parts" lemma
analogous to symmetric balancing in clique-density bounds:

\begin{lemma}[Balancing $|B|=|C|$, plan]
\label{lem:balancing}
Let $W$ be a 3-blowup graphon with parts $A,B,C$ of measures $\alpha,\beta,\gamma$
and a triangle-free filling $F$ inside $A$.  Then there exists another such
graphon $W'$ with $|B'|=|C'|=(\beta+\gamma)/2$, the same $\alpha$, the same
total edge density, and $t(P,W')\le t(P,W)$.
\end{lemma}

Sketch: at fixed $\alpha$, $|B|+|C|=1-\alpha$, the cross-edge counts that go
into $t(P,W)$ depend on $|B|,|C|$ only through the product $|B|\cdot|C|$, which
is maximised at $|B|=|C|$.  The "rooted $K_4$ at an $A$-vertex" count likewise
depends symmetrically on $(|B|,|C|)$ in a way that favours balanced parts.
A careful accounting needs to be done.

\subsection{Replacing the filling by a regular one}

For the filling inside $A$, Jensen's inequality (as in the lemma of §7.2 of the
progress note) shows $D_2(F)=\int d_F(u)^2\dd u\ge q^2$ with equality iff $F$ is $q$-regular.
Since $t(P,W)$ depends on the filling only through $q$ and $D_2(F)$, and
positively on both, the minimum at fixed $\alpha$, fixed $q$ is achieved by
$q$-regular fillings, hence by the standard 2-block construction within $A$.
This is the only piece already fully rigorous.

\section{Plan B: flag-algebra certificates}

A direct flag-algebra computation produces certified lower bounds
$t(P,W) \ge L(x)$ at a finite set of rational $x$.  The available tooling is:
\begin{enumerate}[label=(\arabic*)]
\item set up the SDP relaxation of the flag algebra at fixed degree $d$;
\item solve the SDP numerically;
\item round to a rational certificate, verified exactly.
\end{enumerate}
For Problem~1 the analogous program for $C_5$ at $p=1-1/k$ was carried out by
Bennett--Dudek--Lidický--Pikhurko (2020).  For $C_5$ the construction is
geometric and the certificate at $p=1-1/k$ is straightforward; for our
six-vertex graph $H$ at $p=4/5$, the construction is also circular and the
target $24349/187500$ is a small-denominator rational, so the SDP is likely
tractable.

\begin{problem}[Step B]
\label{prob:B}
Run a degree-7 (i.e.\ over $K_7$-flags) SDP to certify the lower bound
$t(H, W) \ge 24349/187500$ for $W$ with $t(K_2,W) = 4/5$.
\end{problem}

For Problem~2 (the conjecture in this note), one would need certificates
parametric in $x$, which is harder.  A practical hybrid is to compute SDP
certificates at a fine mesh $x_1<\dots<x_n$ in $[2/3,3/4]$ and use a
"convexity-style" interpolation argument.  The mixed curvature established
in §5 of the rigorous certificates note breaks the simplest interpolation
(a single convex certificate cannot work everywhere), so the interpolation
needs to know about the inflection $x^*$.

\section{Plan C: rooted clique-density inequalities}

A more conceptual route is to develop sharp rooted clique-density inequalities
directly, using the form~\eqref{eq:P-integral}.

The quantity $t(P,W) = \int d_W \cdot \kappa_4$ is bounded below by:
\begin{itemize}[itemsep=2pt]
\item the Cauchy-Schwarz upper bound $t(P,W)\le (t(K_2,W)\cdot
t(K_4,W))^{1/2}\cdot\bigl(\E[d_W^2]\bigr)^{1/2}\cdot(\E[\kappa_4^2])^{1/2}$;
this is not useful directly;
\item a rearrangement-inequality lower bound: at fixed marginals
$\mu_d := \text{law of } d_W(X)$ and $\mu_\kappa := \text{law of } \kappa_4(X)$,
$\int d_W \kappa_4$ is minimised by the negative-rearrangement coupling.
\end{itemize}

The substantive question is: what couplings $(d_W, \kappa_4)$ are achievable
by actual graphons under $t(K_2,W)=x$?  The set of achievable couplings is
a region in the $(\mu_d,\mu_\kappa)$ plane (or more generally on the
2-dimensional joint distribution) constrained by both Cauchy-Schwarz and
clique-density inequalities.

\begin{problem}[Step C, qualitative]
\label{prob:C}
Identify the boundary curve of the achievable region in the
$\bigl(\E[d_W^2], \E[\kappa_4]\bigr)$ plane (or a richer one) at fixed
$t(K_2,W)=x$ on $[2/3,3/4]$.  Express $\Psi_P(x)$ as a value of an appropriate
rearrangement-style functional restricted to this boundary.
\end{problem}

This plan is more speculative; it has the advantage that it would also explain
the mixed-curvature phenomenon (the boundary curve itself has a kink, and
$\Psi_P''(x)$ changes sign at that kink).

\section{Concrete next steps}

In order of decreasing tractability:

\begin{enumerate}[label=\textbf{Task \arabic*.},itemsep=4pt]
\item \emph{Pin down the Step A2 calculation} (Lemma~\ref{lem:balancing}):
show that, given a 3-blowup with arbitrary parts, balancing $|B|=|C|$ never
increases $t(P,W)$.  This is a purely algebraic check on a 4-variable
expression and should fit into one page.
\item \emph{Sample a flag-algebra SDP} at a fixed $x_0$ in $(2/3,3/4)$, e.g.\
$x_0 = 17/25$.  The target value is the explicit $\Psi_P(17/25)$ we computed.
Solve in CSDP or SDPA, round, and check.  Outcome: either a clean numerical
SDP that rounds to a rational certificate, or evidence that the SDP at
moderate degree is too loose.
\item \emph{Survey the literature on $K_4$-supersaturation stability} to see
whether Lemma~\ref{lem:K4-stability} (cut-metric stability) is available off
the shelf for the LS 4-partite template.
\item \emph{Treat the analogous $J$ problem in parallel.}  Step~A2 for $J$ is
similar; the rooted-$K_3$-and-$K_4$ identity~\eqref{eq:J-integral} suggests
the same blowup-then-balance argument.
\end{enumerate}

\section{What the mixed-curvature result tells us about the global problem}

The rigorous mixed-curvature result (rigorous-certificates note,
Theorem~5.x) constrains \emph{what kind} of global lower bound is achievable.
In particular:

\begin{itemize}[itemsep=2pt]
\item A single convex/concave certificate cannot prove
$t(P,W)\ge\Psi_P(x)$ uniformly on $[2/3,3/4]$; the certificate must
distinguish $x<x^*_P$ from $x>x^*_P$.
\item A flag-algebra style SDP at a single $x$ that interpolates linearly to
nearby $x$ values would only be valid in a small neighbourhood; the radius of
that neighbourhood is bounded by the distance to $x^*_P$, where the
\emph{rate} of $\Psi_P''$ change is fastest.
\item A symmetrisation-based proof must implicitly distinguish the two
curvature regimes; the most natural way is that the optimal $\alpha=\alpha^*(x)$
itself depends on $x$, and the bifurcation between "$A$ is the small part"
($\alpha<1/3+\varepsilon$) and "$A$ is the larger part" ($\alpha\to 1/2$)
matches the inflection.
\end{itemize}

This observation suggests that Plan A (stability+symmetrisation) is in fact
most natural, because it respects the parametric structure of the minimiser
across the interval.  Plans B and C are more rigid; they will likely need to
be informed by the parametric structure.

\section{Open question: does the global reduction extend beyond $[2/3,3/4]$?}

The reduced curve $\Psi_P$ is defined for $x\in[2/3,3/4]$ by construction
($q$-regular triangle-free filling inside one part of a 3-blowup).  For
$x\in[3/4,4/5]$ (the second Turán interval) the natural analogue would be a
4-blowup with a triangle-free filling inside one part; for $x\in[4/5,5/6]$
a 5-blowup with $K_4$-free filling; etc.  We have not analysed these.

The mixed-curvature result on $[2/3,3/4]$ leaves the structure on
$[3/4,4/5]$ as the next open question.  If the analogous 4-blowup-with-filling
family is the right model there, the analogous reduction analysis would
produce another $\Psi$-function, with its own curvature signature.

\end{document}
